% =============================================================================
% EXAMPLE DOCUMENT
%
% A mock report built on template.tex, used for the preview in the README.
% Topic, text and numbers are invented; the references are real books.
% The preamble is a copy of the template's, with review marks switched off.
%
% Build from this folder:  latexmk -pdf example.tex
% =============================================================================

\documentclass[11pt,a4paper]{article}

% =================================================
% Font: Arial 11 pt or comparable sans-serif.
% helvet is loaded WITHOUT the scaled option on purpose. [scaled] defaults to
% 0.95 and would render every size at 95 % of its declared value: body 10.4 pt
% instead of 11, footnotes and captions 9.5 instead of 10.
% =================================================
\usepackage[T1]{fontenc}
\usepackage{helvet}
\renewcommand{\familydefault}{\sfdefault}
\usepackage{sfmath}   % one font throughout, maths included

% =================================================
% Page Layout: DIN A4, single-sided, 2 cm margins
% =================================================
\usepackage[
  a4paper,
  top=2cm,
  bottom=2cm,
  left=2cm,
  right=2cm,
  footskip=1cm
]{geometry}

% =================================================
% Language, Hyphenation, Justification
% American English, because APA is an American standard and
% biblatex-apa expects the american language mapping below.
% =================================================
\usepackage[american]{babel}
\addto\captionsamerican{\renewcommand{\contentsname}{Table of Contents}}
\usepackage{microtype}
\usepackage{ragged2e}
\justifying
\hyphenpenalty=50   % hyphenation on, as the formatting guidelines require
\tolerance=2000
\emergencystretch=2em

% =================================================
% Line and Paragraph Spacing
% =================================================
\usepackage{setspace}
\onehalfspacing
\setlength{\parindent}{0pt}
\setlength{\parskip}{6pt}

% =================================================
% Page breaking
% No orphans and no widows; 10000 is TeX's "never". \clubpenalty also keeps a
% heading from being followed by a single line before the page turns.
% \raggedbottom lets a page end short instead of stretching the gaps.
% =================================================
\clubpenalty=10000
\widowpenalty=10000
\displaywidowpenalty=10000
\brokenpenalty=10000        % no page break right after a hyphenated line
\raggedbottom

% =================================================
% Packages
% =================================================
\usepackage{graphicx}
\usepackage{amssymb}
\usepackage{tikz}
\usetikzlibrary{positioning,arrows.meta,fit}
\usepackage{pgfplots}
\pgfplotsset{compat=1.18}
\usepackage{listings}
\lstset{
  basicstyle=\ttfamily\fontsize{7}{8.4}\selectfont,
  numbers=left, numberstyle=\ttfamily\fontsize{6}{7}\selectfont\color{black!45},
  numbersep=6pt,
  frame=single, rulecolor=\color{black!35}, framesep=4pt,
  breaklines=true, breakatwhitespace=true, breakindent=12pt,
  showstringspaces=false, tabsize=4,
  keywordstyle=\bfseries, commentstyle=\color{black!55}, stringstyle=\color{black!70},
  captionpos=b, xleftmargin=16pt
}
\usepackage{float}    % [H]: pin a float exactly where it is written
\usepackage{placeins} % \FloatBarrier: keep a float inside the section that cites it
\usepackage{booktabs}
\usepackage{array}
\usepackage{longtable}
\usepackage[table]{xcolor}
\usepackage{arydshln} % dashed rules in a tabular, loaded after array and colortbl
\usepackage{enumitem}
\usepackage[hang,flushmargin]{footmisc}
\usepackage{hyperref}
\usepackage{xurl}     % long URLs break at any character instead of running into the margin
\usepackage{caption}
\usepackage{titlesec}
\usepackage{tocloft}
\usepackage{fancyhdr}
\usepackage{csquotes}

% three hues that stay apart for the common colour vision deficiencies; pair
% each with its own marker and dash pattern so a plot also reads in print
\definecolor{seriesA}{HTML}{2A78D6}
\definecolor{seriesB}{HTML}{EB6834}
\definecolor{seriesC}{HTML}{1BAF7A}

% =================================================
% REVIEW MARKS
% \rev{text}         highlights a span in yellow, e.g. a number not yet measured
% \revn{text}{note}  highlight plus a visible bracketed note
% \revcite{key}      a citation still to be checked; stays when review is off
% \revsug{text}      a suggested reformulation; disappears when review is off
% \revbad{note}      red mark for a statement no source supports
% SET \reviewfalse BEFORE SUBMITTING: it turns every mark off in one place and
% leaves the text itself untouched.
% =================================================
\usepackage{soul}
\newif\ifreview
\reviewfalse                      % off, so the preview shows the final look
\sethlcolor{yellow}
\newcommand{\rev}[1]{\ifreview\hl{#1}\else#1\fi}
\newcommand{\revn}[2]{\ifreview\hl{#1}\,\textcolor{red}{\footnotesize[#2]}\else#1\fi}
\newcommand{\revcite}[1]{\ifreview\hl{\mbox{\parencite{#1}}}\else\parencite{#1}\fi}
\newcommand{\revsug}[1]{\ifreview\hl{[Suggestion: #1]}\else\fi}
\newcommand{\revbad}[1]{\ifreview{\sethlcolor{red!40}\hl{[! #1]}\sethlcolor{yellow}}\else\fi}

% \anchor{id}    a jump destination with no width and no height
% \gl{key}{text} links a term in the body to its glossary or abbreviation entry
\newcommand{\anchor}[1]{\raisebox{0pt}[0pt][0pt]{\hypertarget{#1}{}}}
\newcommand{\gl}[2]{\hyperlink{gl:#1}{#2}}

% =================================================
% APA Citation Setup
% Requires: biber
% =================================================
\usepackage[
  backend=biber,
  style=apa,
  sorting=nyt,
  sortcites=true
]{biblatex}
\DeclareLanguageMapping{american}{american-apa}
\addbibresource{references.bib}
\setlength{\bibhang}{1.27cm}

% Reminder: IU expects a locator on every citation, e.g.
%   \parencite[p.~12]{key}
%   \textcite[pp.~12--14]{key}
%   \parencite[Section~4.4.4]{key}   % non-numeric postnotes print verbatim

% =================================================
% Captions and Tables
% IU labels floats "Fig." and "Tab."; refer to them by number in the text,
% before they appear.
% =================================================
\captionsetup{
  font=small,
  labelfont=bf,
  labelsep=space,
  justification=justified,
  singlelinecheck=false
}
\captionsetup[table]{name=Tab.}
\captionsetup[figure]{name=Fig.}

% Every float needs a source line in 10 pt, own work included:
%   \srcnote{Own representation.}
% \srcline{...} runs the source into a caption that already sits below.
\newcommand{\srcline}[1]{\emph{Source: #1}}
\newcommand{\srcnote}[2][\linewidth]{%
  {\setlength{\parskip}{0pt}\par\nobreak\vspace{6pt}%
   \noindent\parbox{#1}{\raggedright\itshape\fontsize{10}{12}\selectfont Source: #2}\par}}

\renewcommand{\arraystretch}{1.1}
\setlength{\tabcolsep}{8pt}
% inside a float: 9 pt type, tighter rows, less column padding
\newcommand{\tablefontsize}{\fontsize{9}{10.5}\selectfont
  \renewcommand{\arraystretch}{1.0}\setlength{\tabcolsep}{6pt}}

% =================================================
% Heading Depth: max. three levels
% =================================================
\setcounter{secnumdepth}{3}
\setcounter{tocdepth}{3}

% =================================================
% Headings: bold 16/14/11 pt, 12 pt before and 12/6/6 pt after
% =================================================
\titleformat{\section}
  {\normalfont\fontsize{16}{19}\bfseries}
  {\thesection}{0.75em}{}

\titleformat{\subsection}
  {\normalfont\fontsize{14}{17}\bfseries}
  {\thesubsection}{0.75em}{}

\titleformat{\subsubsection}
  {\normalfont\fontsize{11}{13}\bfseries}
  {\thesubsubsection}{0.75em}{}

\titlespacing*{\section}{0pt}{12pt}{12pt}
\titlespacing*{\subsection}{0pt}{12pt}{6pt}
\titlespacing*{\subsubsection}{0pt}{12pt}{6pt}

% A heading only starts where enough room is left below it for a few lines of
% text; otherwise it moves to the next page.
\usepackage{needspace}
\pretocmd{\section}{\Needspace{10\baselineskip}}{}{}
\pretocmd{\subsection}{\Needspace{6\baselineskip}}{}{}

% Floats sit closer to the text around them; the defaults leave visible gaps.
\setlength{\textfloatsep}{12pt plus 2pt minus 2pt}
\setlength{\intextsep}{12pt plus 2pt minus 2pt}
\setlength{\floatsep}{12pt plus 2pt minus 2pt}

% =================================================
% Unnumbered section style
% =================================================
\titleformat{name=\section,numberless}
  {\normalfont\fontsize{16}{19}\bfseries}
  {}{0pt}{}

% =================================================
% Footnotes: sans-serif 10 pt
% =================================================
\renewcommand{\footnotesize}{\fontsize{10}{12}\selectfont}
\setlength{\skip\footins}{20pt plus 4pt minus 2pt}
\setlength{\footnotesep}{10pt}

% =================================================
% Lists
% =================================================
\setlist{
  itemsep=2pt,
  topsep=6pt,
  partopsep=0pt
}

% =================================================
% Table of Contents
% Only level 1 bold. The contents and list titles are forced to the size of an
% unnumbered \section*, so all front-matter headings match.
% =================================================
\renewcommand{\cfttoctitlefont}{\normalfont\fontsize{16}{19}\bfseries}
\renewcommand{\cftloftitlefont}{\normalfont\fontsize{16}{19}\bfseries}
\renewcommand{\cftlottitlefont}{\normalfont\fontsize{16}{19}\bfseries}

\renewcommand{\cftsecfont}{\bfseries}
\renewcommand{\cftsecpagefont}{\bfseries}
\renewcommand{\cftsubsecfont}{\normalfont}
\renewcommand{\cftsubsecpagefont}{\normalfont}
\renewcommand{\cftsubsubsecfont}{\normalfont}
\renewcommand{\cftsubsubsecpagefont}{\normalfont}

\renewcommand{\cftdotsep}{1.5}
\renewcommand{\cftsecleader}{\cftdotfill{\cftdotsep}}
\renewcommand{\cftsubsecleader}{\cftdotfill{\cftdotsep}}
\renewcommand{\cftsubsubsecleader}{\cftdotfill{\cftdotsep}}

% =================================================
% Footer: centred page number, no header
% =================================================
\pagestyle{fancy}
\fancyhf{}
\fancyfoot[C]{\thepage}
\renewcommand{\headrulewidth}{0pt}
\renewcommand{\footrulewidth}{0pt}

% =================================================
% Hyperlinks: black throughout, so the printed paper stays neutral
% =================================================
\hypersetup{
  colorlinks=true,
  linkcolor=black,
  citecolor=black,
  urlcolor=black,
  pdfborder={0 0 0}
}

% =================================================
% Document
% =================================================
% =================================================
% Document
% =================================================
\begin{document}

\pagenumbering{Roman}
\setcounter{page}{1}

\begin{titlepage}
  \thispagestyle{empty}
  \centering
  \vspace*{1.8cm}
  {\fontsize{28}{32}\selectfont\bfseries Sample Report Title\par}
  \vspace{1cm}
  {\Large\textit{A mock document that shows the layout of the template}\par}
  \vspace{3cm}
  \includegraphics[height=2.5cm]{../images/doc_images/iu-logo.jpg}\par
  \vspace{0.8cm}
  {\large Degree Programme: <PROGRAMME>\par}
  \vspace{0.3cm}
  {\normalsize GitHub: \href{https://github.com}{\textcolor[RGB]{51,153,255}{github.com/<USER>/<REPO>}}\par}
  \vfill
  {\fontsize{8pt}{10pt}\selectfont
  \begin{center}
  \begin{tabular}{@{}p{4cm}p{7cm}}
  \textbf{University:} & IU International University of Applied Sciences \\
  \textbf{Course/Module:} & <MODULE CODE> -- <MODULE NAME> \\
  \textbf{Type of Assessment:} & <ASSESSMENT> \\
  \textbf{Name:} & <NAME> \\
  \textbf{Matriculation Number:} & <MATRICULATION NUMBER> \\
  \textbf{Tutor:} & <TUTOR WITH ACADEMIC TITLE> \\
  \textbf{Submission Date:} & <DATE> \\
  \end{tabular}
  \end{center}
  }
\end{titlepage}

\newpage
\setcounter{page}{2}
\tableofcontents

\newpage
\listoftables
\addcontentsline{toc}{section}{List of Tables}

\newpage
\section*{List of Abbreviations}
\addcontentsline{toc}{section}{List of Abbreviations}
{\fontsize{10}{13}\selectfont
\begin{longtable}{@{}p{2.6cm}p{11.4cm}@{}}
  \anchor{gl:apa}APA & American Psychological Association \\
  \anchor{gl:pdf}PDF & Portable Document Format \\
  \anchor{gl:toc}ToC & Table of Contents \\
\end{longtable}}

\newpage
\pagenumbering{arabic}
\setcounter{page}{1}
\setcounter{table}{0}
\setcounter{figure}{0}

% -------------------------------------------------
\section{Introduction}

This document is a mock report. It exists only to show how a paper built with
the template looks. Its topic, its numbers and its results are invented. The
references, in contrast, are real books on typesetting and academic writing,
so that the reference list shows how the American Psychological Association
(\gl{apa}{APA}) style renders actual entries \parencite{apa2020}.

The report follows the structure IU expects from a project report: an
introduction, a main body and a conclusion. The main body describes a
fictional comparison of three layout variants, called A, B and C. Each
variant is a set of typographic settings, applied to the same sample text.
The comparison asks one question: which variant reads best at a fixed page
budget?

Good typography serves the reader, not the writer \parencite[p.~1]{knuth1984}.
The template therefore fixes the formal settings in one place. A writer can
concentrate on the content, while font, spacing and headings already follow
the guidelines \parencite{lamport1994}.

% -------------------------------------------------
\section{Method}

\subsection{Layout Variants}

The three variants differ in line spacing, paragraph gap and the size of the
body text. Tab.~\ref{tab:variants} lists the settings of each variant. Variant
A matches the template. Variants B and C each change one group of settings,
so that every difference between the results can be traced back to a single
cause.

\begin{table}[H]
  \centering
  \caption{Settings of the three layout variants. Variant A is the template
  itself; B and C each change one group of settings.}
  \label{tab:variants}
  {\tablefontsize
  \begin{tabular}{@{}lccc@{}}
    \toprule
    Setting & Variant A & Variant B & Variant C \\
    \midrule
    Body text (pt)        & 11  & 11  & 12  \\
    Line spacing          & 1.5 & 1.2 & 1.5 \\
    Paragraph gap (pt)    & 6   & 6   & 6   \\
    Heading sizes (pt)    & 16/14/11 & 16/14/11 & 18/15/12 \\
    \bottomrule
  \end{tabular}}
  \srcnote{Own representation.}
\end{table}

\subsection{Procedure}

Each variant typesets the same sample text of ten pages. A fictional panel of
readers rates every variant on a scale from 1 to 5. The panel also records
the reading time per page. Long lines and tight spacing are known to slow
readers down \parencite[p.~92]{mittelbach2004}, so the reading time serves as
a second, independent measure next to the rating.

All values in this report are invented for illustration. They follow no
experiment and support no claim.

% -------------------------------------------------
\section{Results}

\subsection{Ratings}

Tab.~\ref{tab:ratings} summarises the invented ratings. Variant A receives
the highest mean rating. Variant B fits the most text on a page, but its
tighter lines lower the rating. Variant C reads comfortably but needs the
most pages for the same text.

\begin{table}[H]
  \centering
  \caption{Mean rating and pages needed for the sample text. Variant A
  balances readability and length best; all values are invented.}
  \label{tab:ratings}
  {\tablefontsize
  \begin{tabular}{@{}lrrr@{}}
    \toprule
    Variant & Mean rating & Std.\ dev. & Pages \\
    \midrule
    \rowcolor{black!8}
    A (template) & 4.2 & 0.6 & 10.0 \\
    B            & 3.4 & 0.9 &  8.6 \\
    C            & 3.9 & 0.7 & 11.8 \\
    \bottomrule
  \end{tabular}}
  \srcnote{Own representation, invented values.}
\end{table}

\subsection{Reading Time}

Fig.~\ref{fig:time} plots the invented reading time over the ten pages of
the sample text. All three variants slow down slightly towards the end.
Variant B stays the slowest throughout. A plot like this one should show the
data and little else \parencite[p.~93]{tufte2001}: no fill, thin rules and a
legend that names every line.

\begin{figure}[H]
  \centering
  \begin{tikzpicture}
    \begin{axis}[
      width=0.72\linewidth, height=5.2cm,
      xlabel={Page of the sample text}, ylabel={Reading time (s)},
      xmin=1, xmax=10, ymin=50, ymax=110,
      xtick={1,...,10},
      tick label style={font=\footnotesize}, label style={font=\footnotesize},
      legend style={font=\footnotesize, draw=none, at={(0.02,0.98)}, anchor=north west},
      axis line style={black!60}, tick style={black!60},
    ]
      \addplot[seriesA, thick, mark=*, mark size=1.5pt]
        coordinates {(1,62)(2,63)(3,63)(4,65)(5,66)(6,66)(7,68)(8,69)(9,70)(10,71)};
      \addplot[seriesB, thick, dashed, mark=square*, mark size=1.5pt]
        coordinates {(1,78)(2,80)(3,82)(4,83)(5,85)(6,88)(7,90)(8,92)(9,95)(10,97)};
      \addplot[seriesC, thick, dotted, mark=triangle*, mark size=2pt]
        coordinates {(1,58)(2,59)(3,61)(4,61)(5,62)(6,64)(7,65)(8,65)(9,67)(10,68)};
      \legend{Variant A, Variant B, Variant C}
    \end{axis}
  \end{tikzpicture}
  \caption{Invented reading time per page for the three variants. Variant B,
  with the tightest line spacing, is the slowest on every page.}
  \label{fig:time}
  \srcnote{Own representation, invented values.}
\end{figure}

\subsection{Summary of Results}

Tab.~\ref{tab:summary} brings both measures together. It ranks the variants
by rating and by reading time. Both rankings agree on the last place.

\begin{table}[H]
  \centering
  \caption{Ranking of the variants by rating and by mean reading time. Both
  measures place variant B last; all values are invented.}
  \label{tab:summary}
  {\tablefontsize
  \begin{tabular}{@{}lrrrr@{}}
    \toprule
    & \multicolumn{2}{c}{Rating} & \multicolumn{2}{c}{Reading time} \\
    \cmidrule(lr){2-3}\cmidrule(l){4-5}
    Variant & Mean & Rank & Mean (s) & Rank \\
    \midrule
    A & 4.2 & 1 & 66.3 & 2 \\
    B & 3.4 & 3 & 87.0 & 3 \\
    C & 3.9 & 2 & 63.0 & 1 \\
    \bottomrule
  \end{tabular}}
  \srcnote{Own representation, invented values.}
\end{table}

% -------------------------------------------------
\section{Conclusion}

The mock comparison favours variant A, the template itself. It reaches the
best rating at a moderate length. Variant C reads slightly faster but costs
almost two extra pages. Under a fixed page budget, that cost decides against
it.

Since every value in this document is invented, the result only demonstrates
the layout. A real report would state its data source, the sample size and
the uncertainty of every number \parencite{strunk2000}. Appendix~A shows how
an extra table looks outside the page budget.

\newpage
\printbibliography[heading=bibintoc,title={References}]

\newpage
\section*{Appendix}
\addcontentsline{toc}{section}{Appendix}

\subsection*{Appendix A: Per-Page Values}\label{app:a}

The table below lists the invented reading times behind Fig.~\ref{fig:time}.

\begin{table}[H]
  \centering
  \caption{Invented reading time per page and variant, in seconds.}
  \label{tab:appendix}
  {\tablefontsize
  \begin{tabular}{@{}lrrrrrrrrrr@{}}
    \toprule
    Page & 1 & 2 & 3 & 4 & 5 & 6 & 7 & 8 & 9 & 10 \\
    \midrule
    Variant A & 62 & 63 & 63 & 65 & 66 & 66 & 68 & 69 & 70 & 71 \\
    Variant B & 78 & 80 & 82 & 83 & 85 & 88 & 90 & 92 & 95 & 97 \\
    Variant C & 58 & 59 & 61 & 61 & 62 & 64 & 65 & 65 & 67 & 68 \\
    \bottomrule
  \end{tabular}}
  \srcnote{Own representation, invented values.}
\end{table}

\end{document}
